% Class-compatibility preamble retained for a reader excerpt of
%   Sahar Diskin, Anna Geisler, "Perfect Matching in Product Graphs and in their
%   Random Subgraphs", arXiv:2404.14020v4, CC BY 4.0.
%
% The source pins its own publisher class (scrartcl); the delivery wrapper accepts
% only the classes article and amsart, so this file supplies the amsart baseline and
% then copies the author's own declaration block mechanically and unchanged
% (pinned-original-source.tex, source lines 30-57: the author's newtheorem
% declarations and the author's macro declarations).  The author's own package
% list of source lines 4-21 is reused only for the entries the retained text needs
% (enumitem, amsmath, amssymb, amsthm, mathtools, xcolor, geometry) and the
% cross-reference packages already needed by the author's own \Cref/\eqref usage
% (hyperref, cleveref); tikz, caption, subcaption, float, graphicx, wrapfig and
% abstract serve figures and the title block that this excerpt does not carry.
% No mathematical macro, statement, symbol or assumption is introduced, changed
% or removed by this file.
\documentclass[10pt,a4paper]{amsart}
\usepackage[shortlabels]{enumitem}
\usepackage{amsmath}
\usepackage{amssymb}
\usepackage{amsthm}
\usepackage{mathtools}
\usepackage{xcolor}
\usepackage[margin=1in]{geometry}
\usepackage{hyperref}
\usepackage{cleveref}
\hypersetup{
  colorlinks=true,
  linkcolor=blue,
  filecolor=magenta,
  urlcolor=cyan
}

\newtheorem{theorem}{Theorem}
\newtheorem{thm}{Theorem}[section]
\newtheorem{corollary}[thm]{Corollary}
\newtheorem{lemma}[thm]{Lemma}
\newtheorem{proposition}[thm]{Proposition}
\newtheorem{claim}[thm]{Claim}
\newtheorem{remark}[thm]{Remark}
\newtheorem{question}[thm]{Question}

\newcommand{\whp}{whp}
\newcommand{\defi}{\operatorname{def}}
\newcommand{\odd}{\operatorname{odd}}
\newcommand{\iso}{\operatorname{iso}}
\newcommand{\cN}{\mathcal{N}}
\newcommand{\cO}{\mathcal{O}}
\newcommand{\cE}{\mathcal{E}}
\newcommand{\cK}{\mathcal{K}}
\newcommand{\cP}{\mathcal{P}}
\newcommand{\cQ}{\mathcal{Q}}

\newcommand{\footremember}[2]{%
  \footnote{#2}%
  \newcounter{#1}%
  \setcounter{#1}{\value{footnote}}%
}

\newcommand{\ag}[1]{\textcolor{purple}{(Anna: #1)}}
\newcommand{\sd}[1]{\textcolor{red}{(Sahar: #1)}}
\title{An exponential upper bound on the number of large sets with a small edge boundary in a product graph}
\author{Sahar Diskin, Anna Geisler}
\date{arxiv:2404.14020v4; CC BY 4.0}
\begin{document}
\maketitle
\noindent\textit{Editorial note.} All delivered mathematical text is the author's own, copied from the pinned author source sources/combinatorics-fl04-vspace-1-5cm-perfect-matching-in-p/pinned-original-source.tex; no statement or proof step was rewritten, reordered or removed.\par\medskip
\noindent\textit{Editorial note.} Preamble provenance: the source pins the publisher class scrartcl, which the delivery wrapper does not accept, so this card uses sources/combinatorics-fl04-vspace-1-5cm-perfect-matching-in-p/retained-preamble.tex, which supplies the amsart baseline and then carries the author's own newtheorem and macro declarations (pinned source lines 30-57) verbatim. No macro was invented and no mathematical symbol was changed.\par\medskip
\noindent\textit{Editorial note.} Bibliography provenance: the source carries no thebibliography environment and ends with bibliographystyle abbrv plus bibliography perc (author source lines 1691-1692), so the bibliography is resolved with BibTeX from the identical copy staged with the same arXiv version under sources/product-graph-matchings-2404.14020v4/perc.bib; no entry was edited and the bibliography style is the author's own abbrv.\par\medskip
\noindent\textit{Editorial note.} The displayed isoperimetric inequalities that the counted proof applies to sub-products are delivered verbatim as a context span, so the label eq:expansion-local used inside the proof resolves in the assembled material; it is the author's own align environment quoted in place.\par\medskip
\noindent\textit{Editorial note.} The label lem:product-connectivity used inside the counted proof is a navigation pointer only; the statement and its proof are delivered by this lane's separate card for that statement, and the identity itself is declared as a named standard prerequisite here.\par\medskip
\noindent\textit{Editorial note.} Retry note: the earlier revision of this card declared a navigation stub for the label eq:bad-expansion-choices, but that label is defined inside the counted proof itself (its displayed bound at source line 727), so the delivered body and the stub defined the same label twice and the build failed with multiply-defined labels. The retry declares no stub at all: the only genuine cross-reference of the counted proof, to the author's lemma lem:product-connectivity, is satisfied by delivering that lemma's statement verbatim as the context span above.\par\medskip
\noindent\textit{Allowed standard prerequisites (editorial).}\begin{itemize}
\item The product structure of the graph on the coordinate subsets, and the isoperimetric inequality on sub-products. The author's displayed inequalities at source lines 223-226, delivered as a context span of this card. Bounds the number of edges a set must cross when it separates a sub-product, which is the mechanism of the counting argument.
\item The relation between the number of coordinates, the dimension parameter d and the order n of the graph. Set up in the author's preliminaries and proof outline (source lines 201-240). Turns the choice m = log\_2(log\textasciicircum{}5 d) into the numerical bound exp(2a/log d).
\end{itemize}\medskip
\begin{lemma}\label{lem:small-boundary-sets}
Let $\frac{n}{d^{\ln^2 d}}\le a\le \frac{n}{2}$. Then, the number of sets $A\subseteq V$ of size $a$ with $e(A,A^c)< a\ln^2d$ is at most $\exp\left\{\frac{2a}{\ln d}\right\}$.
\end{lemma}

\begin{lemma}\label{lem:product-connectivity}
Let $G=\square_{i=1}^t H_i$, where every $H_i$ is connected and $d_i$-regular.
Let $M_i=|V(H_i)|$ and $d=\sum_{i=1}^t d_i$.
Then
\[
  \lambda(G)=
  \min\left\{d,\ \min_{i\in[t]}\lambda(H_i)\prod_{j\ne i}M_j\right\}.
\]
% In particular, if $M_i\ge3$ for every $i \in [t]$, then
% \[
%   \lambda(G)\ge\min\{d,3^{t-1}\}.
% \]
\end{lemma}

\begin{proof}
Let $\mathcal{F}$ be the family of $A\subseteq V$ satisfying the conditions of the lemma.

For $i\in [t]$ and any $A\subseteq V$, let $E_i(A,A^C)\subseteq E(A,A^C)$ be the set of edges in $E(A,A^C)$ corresponding to a change in the $i$-th coordinate, and let $e_{i,A}\coloneqq |E_i(A,A^C)|$. Moreover, given $I\subseteq [t]$, let $e_{I,A}=\sum_{i\in I}e_{i,A}$. We say that $A$ is \textit{bad} with respect to a set of coordinates $I$, if $e_{I,A}< a\ln^2d\cdot \frac{|I|}{t}$. Let $\mathcal{A}_{I}$ be the family of sets $A$ which are bad with respect to some $I\subseteq [t]$. Note that for every fixed $m \in \mathbb{N}$, if $A\in \mathcal{F}$, then there is some $I$ with $|I|=m$ such that $e_{I,A}< a\ln^2d\cdot \frac{|I|}{t}$. Thus,
\[|\mathcal{F}|\le \sum_{\substack{I \subseteq [t]\\|I|=m}} |\mathcal{A}_I| \leq \binom{t}{m} \max_{\substack{I \subseteq [t]\\|I|=m}} |\mathcal{A}_I|.\]
We now set $m=\log_2\left(\ln^5 d\right)$, and turn to estimate $|\mathcal{A}_I|$ for any $I \subseteq [t]$ with $|I|=m$.

For such an $I$ and $v\in V$, let $G(I,v)\coloneqq \square_{i\in [t]\setminus I}\{v_i\} \square_{i\in I}H_i\subseteq G$ be the product graph where we fix the coordinates outside of $I$ to agree with $v$. That is, $G(I, v)$ is a product graph of dimension $|I|=m$. Observe that $2^{|I|}\le|V(G(I,v))|\le M^{|I|}$, and that for every $v\neq u\in V$, $V(G(I,v))$ and $V(G(I,u))$ are either disjoint or identical. Thus, fixing $I$ with $|I|=m$, there are at most $\frac{n}{2^m}$ different subgraphs $G(I,v)$, and their union is $V$. We say that $A$ \textit{intersects non-trivially} with $G(I,v)$ if $V(G(I,v))\cap A\neq \varnothing$ and $V(G(I,v))\setminus A\neq \varnothing$.
Applying \Cref{lem:product-connectivity} to $G(I, v)$, we have that if $A$ intersects non-trivially with $G(I,v)$, then $G(I,v)$ spans at least $d(G(I, V)) \geq |I|=m$ edges of $E(A,A^C)$. Thus, if $A\in \mathcal{A}_I$, we have that $A$ intersects non-trivially at most $\frac{a\ln^2d}{t}$ such subgraphs. Indeed, otherwise, there would be no $I$ with $|I|=m$ and $e_{I,A}<a\ln^2d\cdot \frac{m}{t}$. Therefore, a set $A\in \mathcal{A}_I$ contains at least $$\frac{a-M^m\frac{a\ln^2d}{t}}{M^m} \leq \frac{a}{M^m}$$ subgraphs of the form $G(I, v)$ for $v \in V$, and at most $M^m\frac{a\ln^2d}{t}$ other vertices. We thus obtain that
\begin{align} \label{eq:bad-expansion-choices}
    |\mathcal{A}_I|&\le \binom{\frac{n}{2^m}}{\frac{a}{M^m}}\binom{n}{M^m\frac{a\ln^2d}{t}},
\end{align}
where we first choose from the at most $n/2^m$ many the subgraphs $G(I, v)$ that are contained in $A$ and then the other vertices.
We may bound \eqref{eq:bad-expansion-choices} from above using that $a \geq \frac{n}{d^{\ln^2 d}}$:
\begin{align*}
    &\left(\frac{enM^m}{a}\right)^{\frac{a}{M^m}}\left(\frac{ent}{M^m a\ln^2d}\right)^{M^m\frac{a\ln^2d}{t}}
    \le \left(ed^{\ln^2d}M^m\right)^{\frac{a}{M^m}}\left(\frac{etd^{\ln^2d}}{M^m\ln^2d}\right)^{M^m\frac{a\ln^2d}{t}}\\
    \le&\left(ed^{\ln^2d}\ln^5 d\right)^{\frac{a}{\ln^5d}}\left(\frac{ed^{\ln^2d+1}}{\ln^7d}\right)^{\frac{a\ln^7d}{t}}
    \le \exp\left\{\frac{a}{\ln d}\right\}.
\end{align*}
Altogether, by multiplying with the number of choices for $I$ we obtain that 
\[
|\mathcal{F}|\le \binom{t}{\log_M(\ln^5d)}\exp\left\{\frac{a}{\ln d}\right\}\le \exp\left\{\frac{2a}{\ln d}\right\},
\]
as required.
\end{proof}

\bibliographystyle{plain}
\bibliography{perc}
\end{document}
